\documentclass[sigconf,anonymous,review]{acmart}
\setcopyright{none}
\settopmatter{printacmref=false}
\renewcommand\footnotetextcopyrightpermission[1]{}
\acmConference[AI-Driven Markets '26]{AI-Driven Market Microstructure in Centralized and Decentralized Finance}{November 2026}{Milan, Italy}
\acmYear{2026}
\title{Risk and Turnover Preferences in Learned Portfolio Rebalancing}
\author{Anonymous Author(s)}
\newcommand{\E}{\mathbb{E}}
\newcommand{\Var}{\operatorname{Var}}
\newcommand{\Cov}{\operatorname{Cov}}
\newcommand{\eq}{\mathrm{eq}}
\newtheorem{proposition}{Proposition}
\newtheorem{theorem}{Theorem}
\begin{document}
\begin{abstract}
Risk and turnover penalties jointly determine the trades selected by a portfolio reward. Replacing variance with a second moment changes this trade-off, even when the risk coefficient is subsequently calibrated. For a constrained one-period allocator, we derive the additional turnover penalty introduced by risk-only calibration and show that jointly scaling risk and cost coefficients preserves the mean--variance optimum. The result follows the standard mean auxiliary and is verified against an adapted MVPI solver. A walk-forward study of ten US industry portfolios over 1990--2025 compares induced allocations on common pre-trade states and realised returns on separate strategy paths. Joint calibration recovers the fitted policy across 17 settings; scalar and risk-only adjustments retain discrepancies. Realised performance differences are inconclusive, and finite-budget policy optimisation introduces a larger error than the primary reward-design gap. The study provides an interpretable test of risk--cost preferences in trading models, with monthly proportional costs rather than an order-book or market-impact model.
\end{abstract}
\ccsdesc[500]{Computing methodologies~Reinforcement learning}
\ccsdesc[300]{Applied computing~Economics}
\keywords{portfolio allocation, reward alignment, mean--variance, reinforcement learning, transaction costs}
\maketitle

\section{Introduction}
A trading reward balances the incentive to take risk against the incentive to trade. These incentives are often assigned separate coefficients, yet changing one part of the reward can alter the effective meaning of the other. In particular, replacing a variance penalty with a squared-return penalty and then recalibrating risk need not preserve the intended turnover preference.

This paper studies that interaction in constrained portfolio rebalancing. On each estimated market state, the intended reference is mean return minus covariance risk and proportional turnover cost. We compare it with second-moment rewards, scalar calibrations and a log-minus-variance preference adapted from Jiang et al.\ \cite{jiang2025}. All local comparisons face the same pre-trade portfolio.

An auxiliary-variable identity makes the distortion explicit. Risk-only calibration adds a turnover term whose magnitude depends on the target expected return. Joint risk--cost calibration removes it and agrees with an adapted MVPI method\ \cite{zhang2021}. Historical data then show how this exact fitted relation differs from realised investment performance.

The setting addresses the robustness and interpretation of AI trading objectives. It uses monthly industry portfolios and a first-order proportional cost approximation. It therefore identifies a reward-level source of preference drift before adding execution dynamics; it does not estimate market impact, liquidity provision or order-book behaviour.

\section{Portfolio objectives and validation}
Direct utility optimisation is established in financial learning. Simon et al.\ \cite{simon2026}, Section 4.1, Eq.\ (17), use centred mean--variance preferences, with turnover costs in Section 3.6. Their deep portfolio policy and US stock-characteristic data differ from the common industry-portfolio environment used here. Cong et al.\ \cite{cong2026}, Eqs.\ (19)--(20), optimise trajectory objectives through pathwise gradients. These approaches show that objective-aware learning is feasible.

Other rewards encode different preferences. Jiang et al.\ \cite{jiang2025}, Eq.\ (8), use net log return minus a variance penalty; Eq.\ (19) adds trading restrictions. We adapt that preference using monthly returns, hard simplex constraints and a first-order turnover charge. The comparison changes the reward while holding the estimated state fixed; it does not reproduce their CNN/WaveNet DPG system. Their earlier TD3 study\ \cite{jiang2024}, Eq.\ (4), combines portfolio value, variance and trading costs, motivating explicit treatment of all three terms.

The calibration below follows the mean auxiliary in MVPI\ \cite{zhang2021}. Known-policy validation also has precedents: Chaouki et al.\ \cite{chaouki2020}, Section 3, compare learned and analytical controls, while Gollart and Okhrin\ \cite{gollart2025}, Section 4.2.1 and Table 1, align terminal utility with a DDPG problem. Our contribution is a reproducible test of risk--cost preference drift and a comparison of its magnitude with optimisation error and historical uncertainty.

\section{Constrained alignment with turnover costs}
Let $m$ and $\Sigma\succ0$ be estimated excess-return moments available before trading, $M=\Sigma+mm^\top$, and $\nu=m^\top M^{-1}m$. The feasible set $C$ is the simplex, $p$ is the pre-trade allocation, and $c(u)=\|u-p\|_1$. For risk aversion $\varphi>0$ and cost coefficient $\lambda\ge0$, define
\begin{equation}\label{eq:target}
 F(u)=m^\top u-\varphi u^\top\Sigma u-\lambda c(u),\quad u\in C.
\end{equation}
Its unique optimum $u^*$ is a \emph{fitted reference}, not the unknown true market optimum. The reference specifies a one-period decision under estimated moments.

\begin{proposition}\label{prop:joint}
Let $y^*=m^\top u^*$ and $A=1+2\varphi y^*>0$. The second-moment reward objective
\begin{equation}
 Q(u)=m^\top u-\frac{\varphi}{A}u^\top M u-\frac{\lambda}{A}c(u)
\end{equation}
has the same unique maximiser as $F$.
\end{proposition}
\begin{proof}
For all feasible $u$,
\begin{equation}\label{eq:identity}
 AQ(u)=F(u)-\varphi(m^\top u-y^*)^2+\varphi(y^*)^2.
\end{equation}
At $u^*$ the square vanishes and $F$ is maximal. Since $A>0$, the assertion follows.
\end{proof}
The argument applies to any feasible set with a global maximiser. Convexity and positive definite covariance ensure uniqueness here. The identity follows the mean auxiliary used in MVPI\ \cite{zhang2021}.

If only the risk coefficient is divided by $A$, while the cost coefficient stays $\lambda$, the right-hand side of Eq.\ \eqref{eq:identity} gains the extra term
\begin{equation}\label{eq:costdistortion}
 -2\varphi y^*\lambda c(u).
\end{equation}
When $y^*>0$, this increases the effective cost preference. It need not imply monotone realised turnover on the strategy's evolving path. The unconstrained formula $\varphi(1-\nu)$ likewise need not calibrate a constrained optimum.

Joint calibration uses $u^*$ and therefore serves as a certificate of alignment with a known reference. As a constructive prior-art comparator, we also alternate maximising
\begin{equation}
 (1+2\varphi y)m^\top u-\varphi u^\top Mu-\lambda c(u)
\end{equation}
and setting $y=m^\top u$. This is an adaptation of the MVPI mean auxiliary to the constrained one-period problem. Including the term $-\varphi y^2$ gives a jointly strictly concave objective, since its quadratic block has Schur complement $\Sigma\succ0$. Exact block updates converge to the same target. This provides a constructive comparator with the same optimum.

\section{Historical study}
\subsection{Data and protocol}
We use Kenneth French's ten value-weighted US industry portfolio monthly total returns and the Fama--French risk-free series\ \cite{french}. The archived downloads have a recorded SHA-256 manifest. Estimation data start January 1970; evaluation covers January 1990--December 2025 (432 months). This is a retrospective study on the August 2026 revised vintage, not a point-in-time constituent dataset. Industry portfolios are research portfolios, not directly tradeable funds; their internal reconstruction costs are not included. The provider's CRSP processing and risk-free source changes also form part of this data vintage.

For month $t$, use only the previous 120 months of excess returns to estimate $m$ and the unbiased covariance, shrunk 10\% toward its diagonal. All strategies are long-only and fully invested. The primary settings are $\varphi=5$ and $\lambda=0.001$ (10 basis points per unit of L1 turnover). Pre-trade weights drift with realised gross asset returns. Net portfolio return uses the explicit first-order convention $u_t^\top R_t-\lambda\|u_t-p_t\|_1$, not a market-impact model or exact intramonth execution equation. The experimental protocol fixed these choices before strategy outcomes were inspected.

\subsection{Comparators and measurements}
We solve the covariance-MV target, raw second moment ($\kappa=\varphi$), unconstrained scalar calibration ($\kappa=\varphi(1-\nu)$), reference-based risk-only calibration, joint calibration, and MVPI adaptation. Equal weight and minimum variance provide conventional portfolio controls. Log utility and log-minus-variance adapt the reward preference in Jiang et al.\ \cite{jiang2025}; every method uses the same monthly data and feasible set. The comparison tests induced preferences under common estimates, not which published learning architecture is best.

For alignment, all methods face the target strategy's \emph{same pre-trade state}. We report Euclidean allocation distance and fitted regret $F(u^*)-F(u)$. Each strategy also follows its own pre-trade states for realised performance. This avoids conflating state divergence with the local reward comparison. Quadratic solutions have independent linear-programme first-order certificates; the largest saved common-state certificate is $4.27\times10^{-9}$ in objective units. The observed agreement of calibrated policies is considerably tighter than the acceptance tolerance.

Realised metrics are monthly excess-return MV, annualised excess Sharpe, return, volatility, turnover and drawdown. Uncertainty for paired MV differences uses 1,000 circular moving-block bootstrap resamples, blocks of 12 months, seed 8128. These intervals describe the observed strategy-return sequences, not full estimator-refitting or model-selection uncertainty. Robustness spans 17 prespecified combinations: $\varphi\in\{2,5,10\}$ and cost $\{0,10,25\}$ bp at the primary estimator, plus windows $\{60,120,240\}$ and covariance shrinkage $\{0,0.1,0.5\}$ at the primary preferences. The overlapping primary case is counted once.

\begin{table*}[t]
\centering
\caption{Primary historical comparison, 1990--2025. Distance and fitted regret use common states; realised metrics use each strategy's own path. Gap and MV score are monthly basis-point units. Calibration improves model alignment; the table does not establish superior investment performance.}
\label{tab:historical}
\input{results_table.tex}
\end{table*}

\subsection{Results}
Table~\ref{tab:historical} shows that raw second-moment optimisation has mean allocation distance 0.00843. Unconstrained scalar calibration reduces this to 0.00540; reference-based risk-only calibration retains 0.00510. Joint calibration and MVPI agree with the fitted reference to numerical precision. The adapted log-minus-variance objective has distance 0.01286; log utility alone is much farther away. The differences measure departure from the specified MV preference; log utility represents a different investment objective.

Across all 17 settings, the largest mean joint-calibration distance is $1.54\times10^{-7}$; raw mean distances range from 0.00157 to 0.03682. Thus constraints and trading costs do not break the joint identity, while unconstrained scalar calibration is not a universal repair. Figure~\ref{fig:historical} relates alignment to historical outcomes, and Figure~\ref{fig:robust} gives robustness and uncertainty.

Joint calibration has realised monthly MV lower than raw by 1.315 bp, with 95\% interval $[-3.318,0.343]$ bp. Against scalar calibration the difference is 0.385 bp. Its interval is $[-0.570,1.351]$ bp. Against risk-only calibration the difference is $-0.131$ bp, with interval $[-0.889,0.634]$ bp. All include zero. Minimum variance has both a better realised MV score and higher Sharpe than the fitted target in the primary sample. This illustrates estimation risk: precise implementation of estimated preferences cannot guarantee the best realised allocation. Equal weight has higher Sharpe than the target but lower realised MV, so Sharpe ranking alone does not answer even the ex-post MV comparison.

\begin{figure*}[p]
\centering\includegraphics[width=.94\textwidth]{figures/historical.pdf}
\caption{Left: mean common-state allocation distance; the logarithmic axis exposes numerical agreement for joint calibration and MVPI. Right: realised Sharpe and monthly MV need not rank policies alike. Equal weight has higher Sharpe but lower MV than the fitted target.}
\Description{Allocation distances on the left and realised Sharpe versus mean-variance on the right.}
\label{fig:historical}
\end{figure*}
\begin{figure*}[p]
\centering\includegraphics[width=.94\textwidth]{figures/robustness.pdf}
\caption{Left: alignment across the 17 prespecified settings, ordered by window, risk penalty, cost and shrinkage. Right: paired 12-month-block bootstrap intervals all include zero. Exact fitted alignment does not imply a statistically clear realised MV advantage.}
\Description{Seventeen-setting allocation-distance curves and three paired confidence intervals crossing zero.}
\label{fig:robust}
\end{figure*}

\subsection{Separating optimisation from reward design}
A finite-budget experiment assigns independent softmax policy logits to each of the 432 historical common states. Adam uses learning rate 0.05, initial logits $N(0,0.1^2)$, five matched seeds and 50, 250, 1,000 and 5,000 gradient steps. Each objective uses only its state's lagged estimates. Independent logits isolate optimisation error on each fitted state. They do not test generalisation across states, as a predictive neural policy would.

Let $u_r$ maximise the specified reward and $\widehat u_r$ be the finite-budget allocation. We report the nonnegative own-reward optimisation gap separately from
\begin{align}
 F(u^*)-F(\widehat u_r)
 &= [F(u^*)-F(u_r)]\nonumber\\
 &\quad+[F(u_r)-F(\widehat u_r)].
\end{align}
The second bracket is a \emph{signed} target-value difference; it is not the own-reward optimisation gap. At 5,000 steps, mean own-reward optimisation gaps remain approximately $4\times10^{-6}$, exceeding the exact raw design gap of $1.48\times10^{-6}$. Thus finite-budget error can dominate the reward discrepancy. Figure~\ref{fig:learning} shows why these must be kept separate: improving the raw-reward optimiser does not remove its exact design gap. The target and joint reward have the same limiting optimum but can exhibit different finite-budget trajectories because their objective geometries differ.

\begin{figure*}[p]
\centering\includegraphics[width=.94\textwidth]{figures/learning.pdf}
\caption{Finite-budget softmax optimisation on historical fitted states. Lines are five-seed means, shading the seed range. Left uses each method's own reward; right uses the common MV target. The dotted line is the exact raw reward-design gap. This diagnostic measures optimisation accuracy, not generalisation.}
\Description{Own-reward and target gaps decline with gradient steps; the raw target gap remains above its design floor.}
\label{fig:learning}
\end{figure*}

\section{Implications for trading models}
The risk and turnover coefficients should be interpreted together. Equation~\eqref{eq:costdistortion} identifies the term left behind by calibrating risk alone. When the target expected return is positive, the fitted objective places extra weight on avoiding turnover. This interpretation concerns a fixed pre-trade state; it does not guarantee less trading along a strategy's realised path.

The historical comparison supports the algebraic correction but not a performance ranking in its favour. Raw, scalar and risk-only rewards produce small primary allocation differences, and their paired MV comparisons with joint calibration include zero. Minimum variance performs better than the fitted target in the primary sample. Reward fidelity and the accuracy of moment estimates therefore remain separate concerns.

In execution or market-making models, inventory risk, impact and fill uncertainty introduce further terms. The present result suggests checking the induced trade-off whenever a quadratic reward is rescaled, but it does not establish the corresponding identity for endogenous costs or liquidity. Those extensions require an execution-level model and data. Within the proportional-cost allocation setting, the common-state audit provides a direct, reproducible test of the preference encoded by a reward.

\section*{Acknowledgments}
OpenAI Codex assisted with the research and preparation of this manuscript. Its use covered literature searches, mathematical derivations, code, experiment analysis, drafting and revision. Reported numerical results were computed from the archived experiments; figures were plotted from those outputs.

\clearpage
\bibliographystyle{ACM-Reference-Format}
\bibliography{references}
\end{document}
